\documentclass[a4paper,10pt]{article}
\usepackage[utf8]{inputenc}
\usepackage{graphicx}
\usepackage{subfigure}
\usepackage{fancyhdr}
\usepackage{listings}
\usepackage{ gensymb }


%opening

\pagestyle{fancy} %eigener Seitenstil
\fancyhf{} %alle Kopf- und Fußzeilenfelder bereinigen
\fancyhead[L]{SEP Abgabegespräch SS2015/16} %Kopfzeile links
\fancyhead[C]{} %zentrierte Kopfzeile
\fancyhead[R]{\includegraphics[width=1cm]{../../template/tugraz_bw}} %Kopfzeile rechts
\renewcommand{\headrulewidth}{0.4pt} %obere Trennlinie
\fancyfoot[C]{} %Seitennummer
\renewcommand{\footrulewidth}{0.4pt} %untere Trennlinie


\lstset{language=C++, basicstyle=\ttfamily\tiny}  

\begin{document}

\section*{Matrix}

Schreibe ein einfaches Programm, das eine Matrix einliest und die transponierte Matrix in ein File schreibt. 

Die Klasse soll einen Filenamen übernehmen, in dem die Matrix steht. Mit dem Befehl 'transpose' soll die Matrix transponiert und ins File 'transpose.txt' geschrieben werden. (Transponieren: 1. Zeile einer Matrix wird zur 1. Spalte, 2. Zeile zur 2.Spalte usw.)\\
Wenn das File nicht gültig ist, soll 'invalid file' ausgegeben werden. Nicht gültig ist ein File, wenn nicht jede Zeile die gleiche Anzahl an Elementen hat.

\subsection*{Beispiel}
\begin{verbatim}
int main() {
    Matrix *matrix = new Matrix("matrix.txt");
    matrix->transpose();
    return 0;
}
\end{verbatim}
\subsubsection*{matrix.txt} 
\begin{verbatim}
1 2 3 4
5 6 7 8
9 10 11 12
\end{verbatim}
\subsubsection*{transpose.txt} 
\begin{verbatim}
1 5 9 
2 6 10 
3 7 11 
4 8 12 
\end{verbatim}
\subsection*{Aufruf}
\texttt{./matrix}

\subsection*{Bewertung}
Dokumentation \& Programmierstil,  Dokumentation \& Programmierstil, und Struktur \& Korrektheit (je 1 Punkt)
\newpage


\end{document}
